\BookChapter{3}

\section[第3章\q 永不知足]{第3章\\ 永不知足}

{\kaishu

当富人开始头脑发热……

}

于2019年去世的先锋领航集团（Vanguard）创始人约翰·博格尔（John Bogle）曾讲过一个故事，让我们意识到了关于金钱我们很少思考的一件事：

{\kaishu

在一位富豪在谢尔特岛（Shelter Island）上举办的一次派对上，库尔特·冯内古特\footnote{美国黑色幽默文学的代表人物之一，代表作有《第五号屠宰场》。}（Kurt Vonnegut）告诉他的朋友约瑟夫·海勒\footnote{美国黑色幽默文学的代表人物之一，因《第二十二条军规》一举成名。}（Joseph Heller），他们的东道主——一位对冲基金经理——在某一天内赚到的钱比海勒那部广受欢迎的小说《第二十二条军规》（\textit{Catch-22}）赚到的全部版税还多。海勒回答道：“是的，但是我拥有一样他永远不会有的东西——知足。”

“知足”——我被这个词简单而又强大的力量震撼了。我感受到的震撼来自两个方面：第一，我自己已经在生活中获得了太多；第二，约瑟夫·海勒的话说得再准确不过了。

“知足”是社会中的一个重要因素，但无数人，包括我们当中很多家财万贯和有权有势者，似乎从来都不知道它的含义是什么。

}

海勒的回答充满了智慧，又那么强有力。

让我再来讲两个例子。它们体现了不知足是多么危险，以及我们应该从中学到什么。

顾磊杰（Rajat Gupta）出生于印度的加尔各答，在10多岁时就成了孤儿。人们说起那些起点很高的特权阶层的孩子时，会形容他们“一出生就上了三垒”\footnote{垒球、棒球运动中的最后一个防守位。跑垒员到达三垒后，马上就可以上本垒并得分了。}，但顾磊杰出生时甚至可以说连棒球馆都看不见。

然而，就是从这样的起点开始，他取得了非凡的成就。

40多岁时，顾磊杰就成了世界上最负盛名的咨询公司麦肯锡（McKinsey）的首席执行官。2007年退休后，他继续在联合国和世界经济论坛担任职务。他和比尔·盖茨合作开展过慈善事业。他是5家上市公司的董事会成员。尽管出生在加尔各答的贫民窟，顾磊杰已经成为当今最成功的商业人士之一。

伴随成功而来的是巨额的财富。截至2008年，顾磊杰的身家达到了1亿美元。\footnote{D. McDonald,“Rajat Gupta: Touched by scandal,” Fortune (October 1, 2010).}对大多数人来说，这都是一笔他们终生无法企及的财富。哪怕以5\%的年回报率来计算，这些钱每小时都能产生接近600美元的收益。

他本可以做任何想做的事。

但他想做的不只是一个百万富翁。他想成为一个亿万富翁，而且是朝思暮想。

顾磊杰是高盛集团（Goldman Sachs）的董事会成员。和他共事的是一群世界上最有钱的投资人。一位投资者在谈到私募股权大亨身家的话题时这样描述顾磊杰：“我觉得他想进入那个圈子。那可是个亿万富翁的圈子，对吧？高盛不就是亿万富翁的圈子吗？”\footnote{“Did millionaire Rajat Gupta suffer from billionaire envy?” The Economic Times (March 27, 2011).}

的确如此。于是顾磊杰找到了一个赚钱的法子。

2008年，在高盛面临金融危机的危急时刻，沃伦·巴菲特决定投资50亿美元来拯救高盛。作为高盛董事会成员，顾磊杰在这个消息公布之前就获悉了这笔交易。这个消息价值千金。高盛能否活下来还不确定，但是巴菲特的支援肯定会推动其股价飙升。

在高盛股东大会上获悉这笔交易即将发生的消息16秒后，顾磊杰拨通了一个名叫拉杰·拉贾拉特南（Raj Rajaratnam）的对冲基金经理的电话。电话没有被录音，但在通话结束后拉贾拉特南立即买入了17.5万股高盛股票，因此不难想象顾磊杰在电话里和他进行了怎样的交流。几个小时后，巴菲特和高盛交易的消息被公布，高盛的股价瞬间暴涨。拉贾拉特南在瞬间就赚到了100万美元。

这只是顾磊杰被指控的内幕交易的一例。美国证券交易委员会（SEC）称，顾磊杰的内幕消息共为其公司获得了1,700万美元的收益。

这笔钱赚得很容易，而对检察官们来说，这个案件办起来更容易。

顾磊杰和拉贾拉特南因为内幕交易而双双获刑入狱。他们的职业生涯和名誉也被彻底毁了。

现在再来说说伯尼·麦道夫（Bernie Madoff）吧，他的罪行是众所周知的。麦道夫是继查尔斯·庞兹\footnote{投机商人，“庞氏骗局”的创始人。}（Charles Ponzi）本人之后最臭名昭著的“庞氏骗局”制造者。在其骗局被揭发之前，麦道夫欺骗了投资者20年——讽刺的是，他就是在顾磊杰购买高盛股票的内幕交易发生几周之后阴谋败露的。

麦道夫是一名做市商，其工作就是撮合股票买卖双方达成交易。他深谙其道。1992年，《华尔街日报》（\textit{The Wall Street Journal}）这样描述麦道夫的做市公司：

{\kaishu

他建立了一家高盈利的证券公司——伯尼·麦道夫证券投资公司，其交易额在纽约证券交易所占据了很大比例。麦道夫公司的场外电子交易的日均成交量为7.4亿，相当于纽约证券交易所每天9\%的交易量。麦道夫公司执行交易的速度之快，成本之低，使得它在实际运营中可以从其他经纪公司以每股1美分的价格买入其客户的订单，然后从其中大多数股票的买卖差价中获利。

}

这可不是对一桩尚未被揭发的骗局进行的不准确的描述——麦道夫的做市生意是合法的。麦道夫公司的一名前职员表示，该公司的做市部门每年的利润能达到2,500万～5,000万美元。

伯尼·麦道夫合法的、非欺诈性的业务无论从何种标准看都取得了巨大的成功。这门生意让他非常富有，而且这种富有是在法律框架内实现的。

然而，他依然冒险进行了欺诈。

我们应该问顾磊杰和麦道夫两个人的问题是：为什么在坐拥上亿财富后，他们仍然要冒着搭上全部身家的风险去获得更多的财富？

那些在生存边缘挣扎的人铤而走险的行为是一码事。一名“尼日利亚骗局”\footnote{一种曾经盛行的垃圾邮件诈骗形式。骗子谎称自己是落难的尼日利亚王子，吸引受害者投资。}的惯犯曾告诉《纽约时报》，他为伤害别人而感到内疚，但“贫穷当前，你就感受不到那种痛苦了”。\footnote{J. Nicas,“Facebook Connected Her to a Tattooed Soldier in Iraq. Or So She Thought,” The New York Times (July 28, 2019).}

而顾磊杰和麦道夫的行为则是另一码事了。他们已经拥有了一切——一般人无法想象的财富、名誉、权力和自由——但他们抛弃了这一切，因为他们想要更多。

他们不懂得知足。

这些都是极端的例子。更多人的不知足并没有达到违法犯罪的地步。

一家对冲基金长期资本管理公司（Long-Term Capital Management）的每名员工的个人金融资产都高达数千万至数亿美元，而他们把大部分金融资产都投入了自己负责的基金。然后他们冒着巨大的风险去追求更多财富，结果却在1998年失去了一切——当时正值历史上最大的牛市，也是美国经济发展最强劲的时期。沃伦·巴菲特后来写道：

{\kaishu

为了赚他们并未拥有也不需要的钱，他们拿自己已经拥有并确实需要的东西去冒险了。这是愚蠢至极的做法。冒着失去重要东西的风险去争取并不重要的东西的行为毫无道理可言。

}

以自己拥有和需要的东西为赌注冒险，希望获得自己没有也不需要的东西，这种行为毫无道理。

这个道理浅显易懂，却经常被我们忽视。

我们中很少有人会像顾磊杰或麦道夫那样拥有1亿美元的金融资产，但是本书读者中有相当大的一部分人在人生中某些时候赚到的钱或拥有的财富能够满足生活中所有合理的需求，并能让自己去做想做的许多事。

如果你是其中一员，请记住以下几点。

\textbf{1. 最难的理财技能是让逐利适可而止。}

这是最重要的原则之一。如果目标随着结果水涨船高，那么你就没有理由去追逐更多了，因为你花更大力气获得更多财富后的感受依然和之前一样。当想拥有更多——更多的金钱、权力和声望——的欲望比满足感增长得更快时，事情就变得很危险了。此时，你向前一步，你想追逐的目标就会向前两步。你会感觉自己始终处在落后的位置，而追上目标的唯一办法就是冒更大的险。

现代资本擅长创造两种东西：财富和嫉妒。或许这二者从来都是相辅相成的，正是想要超越其他人的愿望为人们提供了拼命工作的动力。但是，如果不懂得知足，生活中便没有了快乐，正如那句俗话所说，“幸福是你拥有的减去你期待的”。

\textbf{2. 问题的症结在于攀比心态。}

一个棒球界新手一年的工资收入大约是50万美元。不管怎么看，这都是一笔相当可观的收入了，但假设他和迈克·特劳特\footnote{美国著名职业棒球运动员，在2019年“福布斯100名人榜”上排名第55位。}（Mike Trout）是队友，而特劳特签的是一份为期12年、薪酬共计4.3亿美元的合同。二者相比，这位新手的收入简直不值一提。但再看看迈克·特劳特。他折合每年3,600万美元的工资是一笔高得离谱的收入，但如果你想跻身2018年度收入前十名的对冲基金经理，你需要一年赚至少3.4亿美元，\footnote{T. Maloney,“The Best-Paid Hedge Fund Managers Made \$7.7 Billion in 2018,”Bloomberg (February 15, 2019).}而这些人又是特劳特这种收入水平的人会去做比较的。每年能赚3.4亿美元的对冲基金经理做比较的对象则是排名前五的对冲基金经理，而2018年后者的收入至少要达到7.7亿美元。而这几位顶级基金经理要看齐的对象则是像沃伦·巴菲特这样的顶级富豪，后者2018年的个人收入是35亿美元。而巴菲特则可以向杰夫·贝索斯\footnote{购物网站亚马逊（Amazon）创始人，2020年度美国首富。}（Jeff Bezos）看齐，后者2018年的净收入是240亿美元——这笔收入转换成时薪，都比那位在普通人看来已经很富有的棒球界新手一整年赚的还要多。

问题在于，攀比的天花板很高，因此实际上没有人能触摸到它。或者说，攀比就像一场没有人能打赢的战役，取胜的唯一办法是根本不要加入这场战争——用知足的态度接受一切，即使这意味着自己比周围的人逊色。

我的一个朋友每年都会去拉斯维加斯玩。有一年，他问一个发牌员他自己玩什么，在哪个赌场里玩。这位发牌员冰冷而严肃地说：“在拉斯维加斯，唯一能赌赢的方法就是在刚进来的那一刻转身离开。”

总想变得和别人一样富有的游戏的本质也是如此。

\textbf{3. 你在知足心态下获得的东西已经够多了。}

有人觉得，知足的心态是一种保守主义，会让你错过机会和可能性。

我不这样认为。

知足意味着深知它的反义词——永远难以填满的欲望——会将你推向后悔的境地。

如果你想知道自己的饭量到底有多大，那么唯一的方法就是吃到吐为止。很少会有人尝试这么做，因为呕吐的痛苦是任何美食都无法弥补的。但是，出于某些原因，这种认知却没有被迁移到商业和投资行为中。很多人只有到破产或被迫停止时才会停下逐利的脚步。这种情况可能像工作过劳或无法掌控风险投资配置一样，并不是谁故意为之，但也存在像顾磊杰和伯尼·麦道夫那样不计后果地使用不正当的手段攫取每一分钱的人。

不管怎样，如果你无法拒绝潜在的金钱诱惑，那么欲望最终可能将你吞没。

\textbf{4. 世界上很多事都不值得冒险，无论潜在收益有多大。}

在刑满释放后，顾磊杰告诉《纽约时报》，自己通过这件事得到的教训是：

{\kaishu

不要太看重任何东西，不管是名誉、成就还是其他。现在想想，名誉这东西到底哪里重要呢？好吧，我的名誉被毁了，真是太不公平了。但过于在乎名誉只会徒增烦恼。

}

这可能是他能从这段经历中得到的最糟糕的教训了。在我看来，这不过是一个明知自己的名誉已经一去不返却还拼命想挽回它的人进行的无力的开脱。

名誉是无价的。

自由和独立是无价的。

家人和朋友是无价的。

你希望从某些人处获得的敬爱是无价的。

幸福是无价的。

要想留住这些，最好的方法是懂得什么时候停止冒险，以免伤害到它们。懂得在应该停止的时候停止。

好消息是，帮你获得知足心态的强力工具其实非常简单，而且不会让你冒破坏任何东西的风险。它就是我们下一章的主题。
